% ============================================================================
% nbc_2026_appendix.tex
% BACKUP/APPENDIX deck for the invited talk at the Nanyang Blockchain
% Conference 2026 (NTU Singapore, 21--22 August 2026). Companion to
% slides/nbc_2026_core.tex. Sized per output/nbc_pipeline/00_exemplars.md's
% "appendix-as-defense" note (Takeaway 2: budget roughly as much appendix
% material as core content) and framing_1.md Sec. 3 item 7 ("Backup budget:
% roughly matched to core content -- full regression tables for all four P1
% relationships and the full P2 robustness battery ..., the D-on-L
% robustness detail ..., references slide(s) at full density").
%
% Every number on every slide below is traceable to a file under
% output/nbc_pipeline/04_evidence/ -- re-checked directly against the
% underlying .md/.csv, not copied from a narrative synthesis without
% verification. Where that direct check disagrees with a prose summary
% elsewhere in the pipeline, the primary table wins and the disagreement is
% flagged in the frame (see "P1 backup: depth-split detail" below).
%
% Cosmetic language gate applied (docs/nbc-2026-deck-pipeline.md): avoid
% "rather than," "genuinely," "honest"/"honestly," "broader,"
% "deliberate"/"deliberately"; avoid em-dashes; avoid "." as a list
% separator; at most one "X, not Y" construction in this file; no redundant
% trailing "so ..." clause.
%
% NOT YET COMPILED end-to-end with nbc_2026_core.tex as one deck. Standalone
% file, same preamble/theme as core for visual consistency. Next phase:
% latexmk build + visual QA pass alongside the core deck.
% ============================================================================
\documentclass[aspectratio=169]{beamer}

\usepackage[T1]{fontenc}
\usepackage{lmodern}
\usepackage{amsmath}
\usepackage{amssymb}
\usepackage{booktabs}
\usepackage{array}
\usepackage{multirow}
\usepackage{graphicx}
\usepackage{grffile}

% ---------------------------------------------------------------------------
% Same plain, light template as nbc_2026_core.tex. No logo here: branding
% note restricts the UCL CBT logo to the core deck's title slide only.
% ---------------------------------------------------------------------------
\usetheme{default}
\usecolortheme{default}
\setbeamertemplate{navigation symbols}{}
\setbeamertemplate{footline}{%
  \begin{beamercolorbox}[wd=\paperwidth,ht=2.6ex,dp=1.2ex,leftskip=.4cm,rightskip=.4cm]{structure}
    \usebeamerfont{footline}\scriptsize Backup: \insertsectionhead\hfill\insertframenumber\,/\,\inserttotalframenumber
  \end{beamercolorbox}%
}
\setbeamertemplate{itemize items}[circle]
\setbeamerfont{footline}{size=\scriptsize}

\title[Vehicle Currencies in DeFi -- Backup]{The Making of Vehicle Currencies}
\subtitle{Backup: Full Tables, Robustness Detail, and Q\&A Coverage Map}
\author[J.~Xu et al.]{Jiahua Xu \\ {\small with Emre Ozdenoren, Kathy Yuan, and Benjamin Chen}}
\institute{UCL Centre for Blockchain Technologies}
\date{Nanyang Blockchain Conference 2026 \\ NTU Singapore \\ 21--22 August 2026}

\begin{document}

% ============================================================================
% TITLE / DIVIDER
% ============================================================================
\begin{frame}
  \titlepage
\end{frame}

\begin{frame}{What this appendix covers}
  \begin{itemize}
    \item \textbf{Q\&A coverage map:} the questions NTU's blockchain/CS-leaning audience is most likely to ask, each mapped to a specific slide below.
    \item \textbf{P1 backup:} the full lead-lag table (4 relationships $\times$ 4 horizons) and the full $D$-on-$L$ contra-result detail (every robustness split, plus the LINK falsification cells).
    \item \textbf{P2 backup:} the full placebo-date table, the joint-pretrend Wald detail, the volatility split, and the version-construction table.
    \item \textbf{Design notes:} data provenance and cost construction, why V2$\to$V3 and not cross-chain, candidate-set and route construction.
    \item References, at full density.
  \end{itemize}
  {\scriptsize Companion file to \texttt{slides/nbc\_2026\_core.tex}. Every number below is checked directly against a file under \texttt{output/nbc\_pipeline/04\_evidence/}.}
\end{frame}

% ============================================================================
% Q&A COVERAGE MAP
% ============================================================================
\section{Q\&A coverage map}

\begin{frame}{Q\&A coverage map (1/2): mechanics and data}
  \footnotesize
  \begin{center}
  \begin{tabular}{@{}p{5.6cm}p{5.9cm}@{}}
    \toprule
    Likely question & Backup slide \\
    \midrule
    Does $D$ (DirectCostAdvantage) include gas, or only fees/price impact? & ``Data provenance and cost construction'' \\[0.3em]
    Where does the on-chain data come from, and how do you know an indexer/subgraph didn't drop or mis-attribute swaps? & ``Data provenance and cost construction'' \\[0.3em]
    With only 5 tokens, is Driscoll-Kraay / date-clustering actually valid here? & ``P1 backup: full headline table'' (note) \\[0.3em]
    How is an ``indirect route'' identified on-chain? Could a swap count as ``through'' a candidate just because the trader also held it? & ``Candidate set, route construction, and the LINK placebo'' \\[0.3em]
    Could MEV/arbitrage bots mechanically generate the $S$-on-$L$ result instead of liquidity providers doing so? & ``Candidate set, route construction, and the LINK placebo'' \\
    \bottomrule
  \end{tabular}
  \end{center}
\end{frame}

\begin{frame}{Q\&A coverage map (2/2): design choices and the contra-results}
  \footnotesize
  \begin{center}
  \begin{tabular}{@{}p{5.6cm}p{5.9cm}@{}}
    \toprule
    Likely question & Backup slide \\
    \midrule
    Caparros-Chaudhary-Klein (CCK) is cited as the identification template and uses cross-chain gas variation. Why is P2 tested within one chain instead? & ``Why V2$\to$V3: no cross-chain design'' \\[0.3em]
    Why V2$\to$V3 specifically? Why not V1, or a non-Uniswap AMM? & ``Why V2$\to$V3: no cross-chain design'' \\[0.3em]
    Is LINK really a fair placebo, given it's one unit with no fixed effects? & ``P1 backup: LINK falsification detail'' \\[0.3em]
    Why this 5-token candidate set (WETH/USDC/USDT/DAI/WBTC), with no long-tail candidates? & ``Candidate set, route construction, and the LINK placebo'' \\[0.3em]
    Is the V2$\to$V3 null/confound really a design problem, or could it be a coding bug in the panel? & ``P2 backup: joint pretrend detail'' (note) \\
    \bottomrule
  \end{tabular}
  \end{center}
  {\scriptsize Per \texttt{docs/nbc-2026-deck-pipeline.md}: this is the one place NBC's actual audience profile is allowed to shape content, capped and confined to this appendix.}
\end{frame}

% ============================================================================
% P1 BACKUP
% ============================================================================
\section{P1 backup}

\begin{frame}{P1 backup: full headline table, all 4 relationships $\times$ 4 horizons}
  \tiny
  \begin{center}
  \begin{tabular}{@{}l c r r r r r r@{}}
    \toprule
    Relationship & $\tau$ & $\beta$ & SE (cluster-date) & $p$ (cl.-date) & SE (DK) & $p$ (DK) & SE (bootstrap) \\
    \midrule
    $S$ on $L$ ($+$) & 1  & 0.0223 & 0.0015 & $<$0.001 & 0.0015 & $<$0.001 & 0.0024 \\
    $S$ on $L$ ($+$) & 7  & 0.0261 & 0.0016 & $<$0.001 & 0.0023 & $<$0.001 & 0.0030 \\
    $S$ on $L$ ($+$) & 14 & 0.0270 & 0.0017 & $<$0.001 & 0.0030 & $<$0.001 & 0.0032 \\
    $S$ on $L$ ($+$) & 30 & 0.0270 & 0.0016 & $<$0.001 & 0.0036 & $<$0.001 & 0.0040 \\
    \midrule
    $L$ on $S$ ($+$) & 1  & 0.0189 & 0.0091 & 0.039 & 0.0095 & 0.046 & 0.0094 \\
    $L$ on $S$ ($+$) & 7  & 0.0377 & 0.0114 & $<$0.001 & 0.0146 & 0.010 & 0.0130 \\
    $L$ on $S$ ($+$) & 14 & 0.0483 & 0.0148 & 0.001 & 0.0247 & 0.051 & 0.0276 \\
    $L$ on $S$ ($+$) & 30 & 0.0842 & 0.0154 & $<$0.001 & 0.0344 & 0.014 & 0.0350 \\
    \midrule
    $S$ on $D$ ($-$) & 1  & $-$0.0060 & 0.0035 & 0.085 & 0.0036 & 0.095 & 0.0051 \\
    $S$ on $D$ ($-$) & 7  & $-$0.0057 & 0.0036 & 0.117 & 0.0049 & 0.239 & 0.0067 \\
    $S$ on $D$ ($-$) & 14 & $-$0.0122 & 0.0037 & 0.001 & 0.0065 & 0.060 & 0.0084 \\
    $S$ on $D$ ($-$) & 30 & $-$0.0129 & 0.0041 & 0.002 & 0.0084 & 0.123 & 0.0086 \\
    \midrule
    $D$ on $L$ ($-$) & 1  & \textbf{0.0111} & 0.0058 & 0.054 & 0.0054 & 0.039 & 0.0080 \\
    $D$ on $L$ ($-$) & 7  & \textbf{0.0162} & 0.0060 & 0.007 & 0.0085 & 0.057 & 0.0104 \\
    $D$ on $L$ ($-$) & 14 & \textbf{0.0168} & 0.0064 & 0.008 & 0.0113 & 0.139 & 0.0131 \\
    $D$ on $L$ ($-$) & 30 & \textbf{0.0244} & 0.0062 & $<$0.001 & 0.0154 & 0.113 & 0.0156 \\
    \bottomrule
  \end{tabular}
  \end{center}
  \vspace{0.3em}
  \tiny $N=9{,}265$--$9{,}410$; 1{,}853--1{,}882 date clusters, 5 tokens. Bold = wrong-signed vs.\ the predicted sign in parentheses. \textbf{Inference note (Q\&A):} date-clustering (not token-clustering) is intentional: with only 5 tokens, clustering on the token dimension would give 5 clusters, which is invalid asymptotics. Clustering on date ($\sim$1{,}850--1{,}880 clusters) plus Driscoll-Kraay (which runs on $T\to\infty$, $N$ fixed) is the correct tool for $N=5$; independently re-derived to 4 decimal places (Sec.\ 5, \texttt{05\_verification.md}). One nitpick from that check: the finite-sample multiplier uses $k=4$ regressor d.f.\ without an FE adjustment, modestly understating SEs ($\sim$10--12\%), so DK $p=0.046$--$0.060$ above should not be read as robustly significant. Source: \texttt{04\_evidence/p1/p1\_headline\_panel\_results.md}.
\end{frame}

\begin{frame}{P1 backup: $D$-on-$L$ detail --- depth split (and a correction)}
  \scriptsize
  \begin{center}
  \begin{tabular}{@{}l r r r r@{}}
    \toprule
    Depth regime & $\tau$ & $\beta$ & $p$ (DK) & Sign check \\
    \midrule
    High baseline depth & 1  & $-$0.0245 & $<$0.001 & match \\
    High baseline depth & 7  & $-$0.0244 & 0.008 & match \\
    High baseline depth & 14 & $-$0.0372 & 0.002 & match \\
    High baseline depth & 30 & $-$0.0236 & 0.084 & match (n.s.) \\
    Low baseline depth  & 1  & $-$0.0449 & $<$0.001 & match \\
    Low baseline depth  & 7  & $-$0.0541 & $<$0.001 & match \\
    Low baseline depth  & 14 & $-$0.0617 & $<$0.001 & match \\
    Low baseline depth  & 30 & $-$0.0731 & $<$0.001 & match \\
    \bottomrule
  \end{tabular}
  \end{center}
  \vspace{0.1em}
  \begin{alertblock}{Correction against the prose synthesis, checked against the primary table}
    \scriptsize The $D$-on-$L$ wrong sign does \textbf{not} recur here in either regime: all 8 cells above match the predicted negative sign, most significantly. This is the one split where the contra-result reverses back to the theory's prediction, and it should not be read as a fifth confirmation of the wrong-signed pooled result. \texttt{05\_verification.md}'s Sec.\ 3 prose describes this split as ``recurs in the low-depth regime''; the primary table here, checked directly for this appendix, shows the opposite in both regimes. The wrong-signed cells the depth-split tally actually contains come from $S$-on-$D$ and $L$-on-$S$ instead of $D$-on-$L$; see next slides for where the contra-result really does recur.
  \end{alertblock}
  \vspace{-0.3em}
  {\scriptsize Source: \texttt{04\_evidence/p1/p1\_robustness\_depth\_split.md}, cross-checked against \texttt{p1\_robustness\_depth\_split.csv}.}
\end{frame}

\begin{frame}{P1 backup: $D$-on-$L$ detail --- volatility split (recurs in both regimes)}
  \footnotesize
  \begin{center}
  \begin{tabular}{@{}l r r r r@{}}
    \toprule
    Volatility regime & $\tau$ & $\beta$ & $p$ (DK) & Sign check \\
    \midrule
    Calm & 1  & 0.0071 & 0.409 & wrong sign (n.s.) \\
    Calm & 7  & 0.0220 & 0.159 & wrong sign (n.s.) \\
    Calm & 14 & 0.0232 & 0.282 & wrong sign (n.s.) \\
    Calm & 30 & 0.0346 & 0.262 & wrong sign (n.s.) \\
    High & 1  & 0.0117 & 0.106 & wrong sign (n.s.) \\
    High & 7  & 0.0122 & 0.243 & wrong sign (n.s.) \\
    High & 14 & 0.0137 & 0.316 & wrong sign (n.s.) \\
    High & 30 & 0.0186 & 0.219 & wrong sign (n.s.) \\
    \bottomrule
  \end{tabular}
  \end{center}
  \vspace{0.5em}
  Every one of the 8 volatility-split cells for $D$ on $L$ is wrong-signed (positive, predicted negative); none reach DK significance individually, but the consistency of the sign across every horizon and both regimes is itself the finding this split contributes.
  {\scriptsize Source: \texttt{04\_evidence/p1/p1\_robustness\_volatility\_split.csv} (cross-checked against \texttt{p1\_robustness\_battery\_summary.md}).}
\end{frame}

\begin{frame}{P1 backup: $D$-on-$L$ detail --- alternative measures}
  \tiny
  \begin{center}
  \begin{tabular}{@{}l r r r r@{}}
    \toprule
    Variant & $\tau$ & $\beta$ & $p$ (DK) & Sign check \\
    \midrule
    Alt.\ depth (lp\_concentration) & 1  & 0.0517 & 0.013 & wrong sign \\
    Alt.\ depth (lp\_concentration) & 7  & 0.0781 & 0.021 & wrong sign \\
    Alt.\ depth (lp\_concentration) & 14 & 0.0670 & 0.116 & wrong sign (n.s.) \\
    Alt.\ depth (lp\_concentration) & 30 & 0.0923 & 0.095 & wrong sign (n.s.) \\
    Winsorized-mean $D$ & 1  & 0.0133 & $<$0.001 & wrong sign \\
    Winsorized-mean $D$ & 7  & 0.0195 & $<$0.001 & wrong sign \\
    Winsorized-mean $D$ & 14 & 0.0222 & $<$0.001 & wrong sign \\
    Winsorized-mean $D$ & 30 & 0.0286 & $<$0.001 & wrong sign \\
    \$1k common-support $D$ & 1  & 0.0101 & 0.001 & wrong sign \\
    \$1k common-support $D$ & 7  & 0.0116 & 0.004 & wrong sign \\
    \$1k common-support $D$ & 14 & 0.0134 & 0.008 & wrong sign \\
    \$1k common-support $D$ & 30 & 0.0165 & 0.019 & wrong sign \\
    \$100k common-support $D$ & 1  & $-$0.0075 & 0.164 & match (n.s.) \\
    \$100k common-support $D$ & 7  & $-$0.0047 & 0.596 & match (n.s.) \\
    \$100k common-support $D$ & 14 & $-$0.0065 & 0.558 & match (n.s.) \\
    \$100k common-support $D$ & 30 & 0.0035 & 0.802 & wrong sign (n.s.) \\
    \bottomrule
  \end{tabular}
  \end{center}
  \vspace{0.3em}
  \tiny Wrong-signed in 3 of 4 variants at every horizon (12 of 16 cells); the \$100k window is the one variant/horizon combination that turns null instead of wrong-signed, except at $\tau=30$. Source: \texttt{04\_evidence/p1/p1\_robustness\_altmeasure.csv}.
\end{frame}

\begin{frame}{P1 backup: $D$-on-$L$ detail --- period split and full tally}
  \footnotesize
  \begin{center}
  \begin{tabular}{@{}l r r r r@{}}
    \toprule
    Period & $\tau$ & $\beta$ & $p$ (DK) & Sign check \\
    \midrule
    Pre-midpoint  & 1  & 0.0192 & 0.087 & wrong sign (n.s.) \\
    Pre-midpoint  & 7  & 0.0263 & 0.094 & wrong sign (n.s.) \\
    Pre-midpoint  & 14 & 0.0285 & 0.110 & wrong sign (n.s.) \\
    Pre-midpoint  & 30 & 0.0327 & 0.183 & wrong sign (n.s.) \\
    Post-midpoint & 1  & $-$0.0618 & 0.002 & match \\
    Post-midpoint & 7  & $-$0.0717 & 0.014 & match \\
    Post-midpoint & 14 & $-$0.1026 & 0.008 & match \\
    Post-midpoint & 30 & $-$0.1376 & 0.001 & match \\
    \bottomrule
  \end{tabular}
  \end{center}
  \vspace{0.4em}
  \begin{block}{Where the $D$-on-$L$ contra-result actually recurs, checked directly}
    \footnotesize Pooled headline (wrong-signed, all 4 $\tau$) $\to$ \textbf{volatility split: recurs in both regimes} $\to$ \textbf{alt-measure: recurs in 3 of 4 variants} $\to$ \textbf{period split: recurs pre-midpoint, reverses post-midpoint} $\to$ \textbf{depth split: reverses in both regimes} (previous slides). Net: this contra-result survives most but not all of the battery. A "nearly every split" characterization does not hold on its own; two of five checks (depth, post-midpoint period) actually flip back to the predicted sign.
  \end{block}
  {\scriptsize Source: \texttt{04\_evidence/p1/p1\_robustness\_period\_split.csv}.}
\end{frame}

\begin{frame}{P1 backup: LINK falsification detail}
  \footnotesize
  \begin{center}
  \begin{tabular}{@{}l r r r r r@{}}
    \toprule
    Equation & $\tau$ & $\beta$ (LINK) & $p$ (NW, LINK) & $\beta$ (pooled) & Same sign? \\
    \midrule
    $S$ on $L$ & 1  & 0.0003  & 0.380 & 0.0223 & yes \\
    $S$ on $L$ & 7  & 0.0001  & 0.911 & 0.0261 & yes \\
    $S$ on $L$ & 14 & $-$0.0008 & 0.183 & 0.0270 & no \\
    $S$ on $L$ & 30 & 0.0001  & 0.898 & 0.0270 & yes \\
    $L$ on $S$ & 1  & $-$0.6536 & 0.152 & 0.0189 & no \\
    $L$ on $S$ & 7  & $-$2.4480 & 0.036 & 0.0377 & no (wrong sign, significant) \\
    $L$ on $S$ & 14 & $-$1.4347 & 0.363 & 0.0483 & no \\
    $L$ on $S$ & 30 & 0.1612  & 0.900 & 0.0842 & yes \\
    \bottomrule
  \end{tabular}
  \end{center}
  \vspace{0.3em}
  \scriptsize\textbf{Is LINK a fair placebo (Q\&A)?} Liquid, real V3 pools, a genuine occasional intermediate hop. It is not conventionally treated as a vehicle/bridge currency. Built by the identical $L$/$S$ construction. But: single cross-sectional unit, no token+date FE possible, Newey-West HAC in place of Driscoll-Kraay -- a strictly weaker, noisier estimator. Sign matches the pooled headline in only 4 of 8 cells; LINK's own coefficients are indistinguishable from zero in 7 of 8; one cell is wrong-signed and significant ($p=0.036$). No placebo $D$ was built: doing so would require rerunning the multi-year V2+V3 quote-simulation with LINK added, out of scope for this pass. This limit is disclosed in \texttt{build\_link\_placebo\_panel.py}; it was not silently skipped. \textbf{Verdict: a weak, inconclusive placebo. It is not a clean pass either way.}

  {\scriptsize Source: \texttt{04\_evidence/p1/p1\_robustness\_placebo\_link.md}.}
\end{frame}

\begin{frame}{P1 backup: full battery, remaining cells}
  \begin{center}
  \begin{block}{TODO -- formatting pass}
    \small The subsample sign-check tallies (depth 72 cells, volatility 72, alt-measure 144, period 72 -- 360 cells total) are summarized on this appendix's earlier slides for the $D$-on-$L$ relationship specifically. The remaining cells (the $S$-on-$L$, $L$-on-$S$, and $S$-on-$D$ relationships across every split) are all directionally consistent with the headline (see \texttt{p1\_robustness\_battery\_summary.md}'s tallies: $S$-on-$L$ is significant in all 32/32 depth$\times$volatility$\times$period$\times$horizon cells and all 16/16 alt-measure$\times$horizon cells). A full row-by-row table across all four relationships and all four split types is not reproduced here; pull directly from the CSVs below if a specific cell is requested live.
  \end{block}
  \end{center}
  {\scriptsize Sources: \texttt{04\_evidence/p1/p1\_robustness\_\{depth\_split,volatility\_split,altmeasure,period\_split\}.csv}.}
\end{frame}

% ============================================================================
% P2 BACKUP
% ============================================================================
\section{P2 backup}

\begin{frame}{P2 backup: placebo dates, level DiD}
  \footnotesize
  \begin{center}
  \begin{tabular}{@{}l l r r@{}}
    \toprule
    Placebo date (regime) & Outcome & Effect & $p$ \\
    \midrule
    2020-10-01 (pre) & Route share (pp) & $-$1.7776 & 0.049 \\
    2020-10-01 (pre) & Vehicle HHI & $-$0.0224 & 0.049 \\
    2020-10-01 (pre) & Direct-pool depth & 0.0554 & 0.003 \\
    2020-11-15 (pre) & Route share (pp) & $-$0.0901 & 0.873 \\
    2020-11-15 (pre) & Vehicle HHI & $-$0.0261 & 0.002 \\
    2020-11-15 (pre) & Direct-pool depth & 0.0211 & 0.048 \\
    2021-01-01 (pre) & Route share (pp) & 2.0903 & $<$0.001 \\
    2021-01-01 (pre) & Vehicle HHI & 0.0018 & 0.768 \\
    2021-01-01 (pre) & Direct-pool depth & 0.0125 & 0.077 \\
    2022-05-05 (post, +1y) & Route share (pp) & $-$2.4767 & $<$0.001 \\
    2022-05-05 (post, +1y) & Vehicle HHI & $-$0.0287 & $<$0.001 \\
    2022-05-05 (post, +1y) & Direct-pool depth & $-$0.0255 & 0.043 \\
    2023-05-05 (post, +2y) & Route share (pp) & 2.9108 & $<$0.001 \\
    2023-05-05 (post, +2y) & Vehicle HHI & 0.0150 & 0.061 \\
    2023-05-05 (post, +2y) & Direct-pool depth & $-$0.0067 & 0.301 \\
    \bottomrule
  \end{tabular}
  \end{center}
  \vspace{-0.3em}
  {\scriptsize Level DiD, $\pm$120-day windows entirely inside one regime (never straddling 2021-05-05). Source: \texttt{04\_evidence/p2/p2\_robustness\_battery\_results.md}, part (i).}
\end{frame}

\begin{frame}{P2 backup: placebo dates, elasticity test}
  \footnotesize
  \begin{center}
  \begin{tabular}{@{}l l r r@{}}
    \toprule
    Placebo date (regime) & Outcome ($|$day-to-day $\Delta|$) & Effect & $p$ \\
    \midrule
    2020-10-01 (pre) & Route share (pp) & 3.4318 & $<$0.001 \\
    2020-10-01 (pre) & Vehicle HHI & 0.0148 & 0.004 \\
    2020-10-01 (pre) & Direct-pool depth & $-$0.0192 & $<$0.001 \\
    2020-11-15 (pre) & Route share (pp) & 4.0151 & $<$0.001 \\
    2020-11-15 (pre) & Vehicle HHI & 0.0207 & $<$0.001 \\
    2020-11-15 (pre) & Direct-pool depth & $-$0.0140 & 0.001 \\
    2021-01-01 (pre) & Route share (pp) & 5.9643 & $<$0.001 \\
    2021-01-01 (pre) & Vehicle HHI & 0.0187 & $<$0.001 \\
    2021-01-01 (pre) & Direct-pool depth & $-$0.0098 & 0.003 \\
    2021-05-05 (\textbf{true launch}) & Route share (pp) & $-$0.6810 & 0.621 \\
    2022-05-05 (post, +1y) & Route share (pp) & $-$0.6550 & 0.427 \\
    2022-05-05 (post, +1y) & Direct-pool depth & 0.0219 & $<$0.001 \\
    2023-05-05 (post, +2y) & Route share (pp) & 1.1534 & 0.135 \\
    2023-05-05 (post, +2y) & Direct-pool depth & $-$0.0101 & 0.036 \\
    \bottomrule
  \end{tabular}
  \end{center}
  \vspace{-0.3em}
  \footnotesize Every pre-launch placebo produces a large, significant ``elasticity'' jump; the true launch window does not. This is the falsification check failing at the pre-launch placebos and the pattern that motivates the pretrend test on the next slide.
  {\scriptsize Source: \texttt{04\_evidence/p2/p2\_robustness\_battery\_results.md}, part (i).}
\end{frame}

\begin{frame}{P2 backup: joint pretrend detail}
  \footnotesize
  \begin{center}
  \begin{tabular}{@{}l l r r r r@{}}
    \toprule
    Window & Outcome & Slope & Slope $p$ & Joint $\chi^2$ (12 df) & Joint $p$ \\
    \midrule
    $\pm$12mo & Route share & 0.0012 & 0.994 & 64.24 & $<$0.001 \\
    $\pm$12mo & Vehicle HHI & $-$0.0024 & 0.144 & 50.52 & $<$0.001 \\
    $\pm$12mo & Direct-pool depth & 0.0176 & $<$0.001 & 60.33 & $<$0.001 \\
    $\pm$24mo & Route share & 0.0344 & 0.836 & 67.48 & $<$0.001 \\
    $\pm$24mo & Vehicle HHI & $-$0.0020 & 0.202 & 49.35 & $<$0.001 \\
    $\pm$24mo & Direct-pool depth & 0.0184 & $<$0.001 & 67.51 & $<$0.001 \\
    \bottomrule
  \end{tabular}
  \end{center}
  \vspace{0.6em}
  A flat \emph{linear} pretrend slope alongside a hugely significant \emph{joint} test that all pre-launch month dummies are zero is a step-like, non-monotonic calendar effect. It is not a smooth trend toward the launch -- exactly the signature the raw monthly $|\Delta$ route share$|$ series shows (16$\to$21pp step-up, Dec 2020$\to$Jan 2021, peaking at launch month, falling back by June 2021).
  \vspace{0.4em}
  \textbf{Bug or design problem (Q\&A)?} \texttt{05\_verification.md} traced this with an independent code audit (the placebo estimator and headline estimator import and call the identical functions, so a latent bug would bias both in the same direction instead of selectively failing only the placebos) plus a raw-panel recomputation bypassing both scripts. Verdict: 85\% confidence this is a genuine design/confound issue and not a coding error.
  {\scriptsize Source: \texttt{04\_evidence/p2/p2\_robustness\_battery\_results.md}, part (i), joint pretrend block.}
\end{frame}

\begin{frame}{P2 backup: volatility split, both windows}
  \footnotesize
  \begin{center}
  \begin{tabular}{@{}l l l r r@{}}
    \toprule
    Window & Pre-V3 volatility & Outcome ($|\Delta|$) & Effect & $p$ \\
    \midrule
    $\pm$12mo & High & Route share (pp) & 0.4028 & 0.823 \\
    $\pm$12mo & Low  & Route share (pp) & $-$1.1501 & 0.522 \\
    $\pm$12mo & High & Vehicle HHI & $-$0.0030 & 0.778 \\
    $\pm$12mo & High & Direct-pool depth & $-$0.0163 & 0.102 \\
    $\pm$12mo & Low  & Vehicle HHI & $-$0.0044 & 0.579 \\
    $\pm$12mo & Low  & Direct-pool depth & 0.0047 & 0.137 \\
    $\pm$24mo & High & Route share (pp) & 1.6682 & 0.295 \\
    $\pm$24mo & Low  & Route share (pp) & $-$0.9029 & 0.670 \\
    $\pm$24mo & High & Vehicle HHI & 0.0010 & 0.906 \\
    $\pm$24mo & High & Direct-pool depth & $-$0.0120 & 0.278 \\
    $\pm$24mo & Low  & Vehicle HHI & $-$0.0052 & 0.530 \\
    $\pm$24mo & Low  & Direct-pool depth & 0.0143 & 0.002 \\
    \bottomrule
  \end{tabular}
  \end{center}
  \vspace{0.4em}
  Null in both volatility regimes for route share and vehicle HHI, every window. The one significant cell (low pre-V3 volatility, $\pm$24mo, direct-pool depth, $p=0.002$) is a capital-efficiency result. It is distinct from a vehicle-status/elasticity result, consistent with the headline's separation of the two margins.
  {\scriptsize Source: \texttt{04\_evidence/p2/p2\_robustness\_battery\_results.md}, part (ii).}
\end{frame}

\begin{frame}{P2 backup: version construction, $\pm$12mo}
  \tiny
  \begin{center}
  \begin{tabular}{@{}l r r r@{}}
    \toprule
    Version construction & Outcome & Effect & $p$ \\
    \midrule
    Best-across-versions & Direct-route availability (pp) & 11.13 & $<$0.001 \\
    Best-across-versions & Vehicle-route availability (pp) & 25.99 & $<$0.001 \\
    Best-across-versions & No-direct vehicle-route availability (pp) & $-$2.16 & 0.013 \\
    Best-across-versions & Direct-route quality (ratio) & 0.139 & $<$0.001 \\
    Best-across-versions & Direct cost advantage vs.\ vehicle route & $-$0.012 & 0.663 \\
    V2-only & Direct-route availability (pp) & $-$39.66 & $<$0.001 \\
    V2-only & Vehicle-route availability (pp) & $-$28.18 & $<$0.001 \\
    V2-only & No-direct vehicle-route availability (pp) & $-$2.29 & $<$0.001 \\
    V2-only & Direct-route quality (ratio) & 0.024 & $<$0.001 \\
    V2-only & Direct cost advantage vs.\ vehicle route & $-$0.005 & 0.821 \\
    V3-only & Direct-route availability (pp) & 50.79 & $<$0.001 \\
    V3-only & Vehicle-route availability (pp) & 28.77 & $<$0.001 \\
    V3-only & No-direct vehicle-route availability (pp) & 8.52 & $<$0.001 \\
    V3-only & Direct-route quality (ratio) & $-$0.0002 & 0.321 \\
    V3-only & Direct cost advantage vs.\ vehicle route & 0.0000 & 0.574 \\
    \bottomrule
  \end{tabular}
  \end{center}
  \vspace{0.3em}
  \tiny ``Best-across-versions'' is the default unfiltered panel (whichever version gave the best realized quote per leg); V2-only/V3-only restrict to legs whose realized best source was actually that version -- a disclosed proxy for version availability. It is not a forced re-quote. Availability shifts mechanically with which version existed; cost-advantage effects are small and insignificant throughout. Source: \texttt{04\_evidence/p2/p2\_robustness\_battery\_results.md}, part (iii).
\end{frame}

\begin{frame}{P2 backup: version construction, $\pm$24mo}
  \tiny
  \begin{center}
  \begin{tabular}{@{}l r r r@{}}
    \toprule
    Version construction & Outcome & Effect & $p$ \\
    \midrule
    Best-across-versions & Direct-route availability (pp) & 11.51 & $<$0.001 \\
    Best-across-versions & Vehicle-route availability (pp) & 26.07 & $<$0.001 \\
    Best-across-versions & No-direct vehicle-route availability (pp) & $-$2.31 & 0.012 \\
    Best-across-versions & Direct-route quality (ratio) & 0.129 & $<$0.001 \\
    Best-across-versions & Direct cost advantage vs.\ vehicle route & $-$0.002 & 0.932 \\
    V2-only & Direct-route availability (pp) & $-$43.89 & $<$0.001 \\
    V2-only & Vehicle-route availability (pp) & $-$29.16 & $<$0.001 \\
    V2-only & No-direct vehicle-route availability (pp) & $-$2.91 & $<$0.001 \\
    V2-only & Direct-route quality (ratio) & 0.017 & 0.003 \\
    V2-only & Direct cost advantage vs.\ vehicle route & 0.001 & 0.970 \\
    V3-only & Direct-route availability (pp) & 55.40 & $<$0.001 \\
    V3-only & Vehicle-route availability (pp) & 32.22 & $<$0.001 \\
    V3-only & No-direct vehicle-route availability (pp) & 7.95 & $<$0.001 \\
    V3-only & Direct-route quality (ratio) & $-$0.0002 & 0.320 \\
    V3-only & Direct cost advantage vs.\ vehicle route & 0.0000 & 0.873 \\
    \bottomrule
  \end{tabular}
  \end{center}
  \vspace{0.3em}
  \tiny Same pattern as $\pm$12mo, both magnitude and significance. Source: \texttt{04\_evidence/p2/p2\_robustness\_battery\_results.md}, part (iii).
\end{frame}

\begin{frame}{P2 backup: fixed-window reproduction (robustness item iv)}
  \footnotesize
  The level-DiD and elasticity-test numbers under a fixed 12- and 24-month event window reproduce the headline exactly (Effect/SE/$t$/$p$ identical to the headline table in \texttt{p2\_vehicle\_status\_elasticity\_results.md}, to the reported decimal places). This confirms the headline windows are not a hand-picked variant that happens to look best: they are the same fixed-window construction the robustness battery independently re-derives.
  {\scriptsize Source: \texttt{04\_evidence/p2/p2\_robustness\_battery\_results.md}, part (iv).}
\end{frame}

% ============================================================================
% DESIGN / DATA NOTES (answers the mechanics Q&A)
% ============================================================================
\section{Design notes}

\begin{frame}{Data provenance and cost construction}
  \scriptsize
  \begin{itemize}
    \item \textbf{Source:} swap/mint/burn/pool-day streams fetched from The Graph subgraphs (six sources: Uniswap V2/V3/V4, Curve, SushiSwap V3, Balancer, each pinned to a specific verified deployment ID) plus Dune for one venue (Fluid); stored as verbatim gzipped JSONL with per-day block-provenance sidecars (min/max block, head-block-at-fetch).
    \item \textbf{Reliability check:} a drift spot-audit re-fetches a random sample of days per source and diffs row IDs, counts, and block ranges against the stored archive. A drifting source is refetched in full; it is not patched cell by cell.
    \item \textbf{USD repricing:} subgraph-native \texttt{amountUSD}/\texttt{amountInUSD} fields are not trusted directly (a known \texttt{derivedETH} blow-up issue in the upstream derivation). Every swap leg is repriced from raw token amounts against a stablecoin-anchored per-day price table instead.
    \item \textbf{DirectCostAdvantage ($D$), does it include gas?} Yes. It is an all-in measure: registered pool fee, price impact from quote replay, and historical route-specific gas expenditure (route gas usage $\times$ day-$t$ gas-token USD price), computed identically for the direct and indirect leg and differenced. Fee/impact/gas components are retained separately and sum back to the all-in figure within a numerical tolerance.
    \item \textbf{Candidate liquidity ($L_{k,t}$):} sum of Uniswap V3 pool TVL where the candidate is a pool side, restricted to pools with $0<\mathrm{TVL}\le\$10$bn (a spam-pool filter) and confirmed token contracts in the persisted V3 swap archive.
  \end{itemize}
  {\scriptsize Sources: \texttt{docs/data-reuse-plan.md}; \texttt{docs/research-questions-and-empirical-design.md} (RQ1 Experiment A; variable registry entries for $C^{D,\mathrm{gas}}$, $C^{I,\mathrm{gas}}$); \texttt{src/ddvc/variable\_registry.py}.}
\end{frame}

\begin{frame}{Why V2$\to$V3: no cross-chain design}
  \scriptsize
  \begin{itemize}
    \item Caparros, Chaudhary \& Klein (CCK, 2024) is cited only for its own result (gas costs causally raise LP repositioning and concentration, across chains) as the identification template P2 was originally scoped to reuse.
    \item That specific reuse was scoped and found not closeable in this pass: this repo's data provenance is Ethereum-mainnet subgraph deployments only, with no multi-chain or L2 on-chain fetch path. Building one is real data-engineering work with its own new provenance and reliability burden. It is not a rerun of an existing pipeline.
    \item The V2$\to$V3 concentrated-liquidity launch (2021-05-05) is the closest available within-chain shock to LP repositioning cost. Chronological sample splits ($\pm$12mo/$\pm$24mo, placebo dates, pretrend tests) substitute for the cross-chain contrast as this design's robustness axis. Java confirmed this substitution up front as the intended design.
    \item Cost of the substitution: it collapses CCK's actual source of variation (cross-chain gas-cost heterogeneity) into a single across-time architecture change, which is exactly why a market-wide volatility episode centered on the same calendar date becomes a live, undissolvable confound for the elasticity outcome (previous section). A cross-chain replication remains future work; it is not claimed as done here.
    \item Why V2/V3 and not V1: this repo's tick-level concentrated-liquidity construction is V3-only and populated from 2021-05-04 onward; extending it to V1/V2-equivalent liquidity is the same kind of non-trivial build work already declined for P1's exogenous-shock search (the September 2020 UNI liquidity-mining launch predates the V3-only $L_{k,t}$ construction for the identical reason).
  \end{itemize}
  {\scriptsize Sources: \texttt{02\_framings/framing\_1.md} Sec.\ 1 (P2 revision); \texttt{docs/data-reuse-plan.md}.}
\end{frame}

\begin{frame}{Candidate set, route construction, and the LINK placebo}
  \scriptsize
  \begin{itemize}
    \item \textbf{Candidate set} $\mathcal K=\{\mathrm{WETH},\mathrm{USDC},\mathrm{USDT},\mathrm{DAI},\mathrm{WBTC}\}$: prespecified in \texttt{src/ddvc/variable\_registry.py}, the largest-TVL, most consistently bridging tokens across the sampled AMMs over the full window. It is not re-optimized on the outcome. A no-long-tail candidate set is a real scope limit, disclosed here: it has not been tested against smaller/altcoin candidates in this pass.
    \item \textbf{How ``indirect route'' is identified:} realized swap paths are reconstructed from the raw swap archive and classified by route class (direct single-hop vs.\ a coherent multi-hop path via a candidate). VehicleShare/BridgeShare counts a route as ``through $k$'' only when $k$ is an actual intermediate hop of a realized multi-leg swap. Merely holding the token alongside the trade does not count.
    \item \textbf{MEV/arbitrage bots:} this construction does not separately identify bot-driven multi-hop arbitrage flow from other trader flow. Bot flow is not excluded. The lead-lag design's own stated limit covers this case: it establishes temporal precedence only, and does not rule out a common shock, such as a coordinated bot response to deepening liquidity, as an alternative to LP-driven feedback.
    \item \textbf{LINK as the falsification token:} chosen because it is liquid, has real V3 pools, and is a genuine occasional intermediate hop, but is not conventionally treated as a vehicle/bridge currency. See the dedicated LINK slide (P1 backup) for why the result reads as weak and inconclusive. It is not presented as a clean pass.
  \end{itemize}
  {\scriptsize Sources: \texttt{src/ddvc/variable\_registry.py}; \texttt{04\_evidence/p1/build\_link\_placebo\_panel.py} (docstring); \texttt{docs/research-questions-and-empirical-design.md} RQ1.}
\end{frame}

% ============================================================================
% REFERENCES
% ============================================================================
\section{References}

\begin{frame}{References}
  \scriptsize
  \begin{itemize}
    \item Somogyi, F.\ (2026). FX vehicle-currency routing, quasi-experimental holiday design. Motivation only.
    \item Krugman, P.\ (1980). Three-country transaction-cost vehicle-selection mechanism. Motivating analogy for P1.
    \item Yuan, K.\ (2005). Benchmark/common-factor liquidity, conditional theory. Motivating analogy for P1.
    \item Leh\'ar, A.\ and Parlour, C.\ (2024). AMM (V1/V2) pool size and price impact vs.\ a limit order book. Motivating analogy for P1.
    \item Klein, O., Kozhan, R., Viswanath-Natraj, G., and Wang, J.\ (2026). LP minting/burning and short-horizon return predictability. Motivates strategic LP behavior behind P1.
    \item Caparros, B., Chaudhary, S., and Klein, O.\ (2024). Gas costs, LP repositioning, and liquidity concentration. Identification template for P2 (cross-chain design not reused here; see ``Why V2$\to$V3: no cross-chain design'').
    \item Klein, O.\ and Song, Y.\ (2021); He, Z., Khorrami, P., and Song, Y.\ (2022). TradFi liquidity-commonality and intermediary-capacity analogies. Background only, no DeFi extension claimed.
    \item Barbon, A.\ and Ranaldo, A.\ (2024/Articles in Advance). CEX vs.\ DEX market quality, gas-fee causal effects. DeFi-native existence-proof comparator for venue/scope calibration. It is not a design template used above.
  \end{itemize}
  {\scriptsize Full corpus manifest: \texttt{output/nbc\_pipeline/00\_manifest.md}. This appendix's own sources: \texttt{output/nbc\_pipeline/04\_evidence/\{p1,p2\}/}, \texttt{docs/research-questions-and-empirical-design.md}, \texttt{docs/data-reuse-plan.md}.}
\end{frame}

\end{document}
